\section{Finite interpolation}\label{sec:interpolation}

The depth-first and entropy estimates let us move two large distinct
clusters from the ends of a path to the ends of one edge. Binary
subdivision limits the number of new clusters used to a logarithm of
the path length. We must also control the probability cost of joining
these clusters; counting the opened edges alone does not control how
many original configurations can produce the same output.

\begin{proposition}[Interpolation along a path]\label{prop:interpolation}
Let $\mu$ be Bernoulli bond percolation with parameter $p\in(0,1)$ on
a finite undirected graph, and put $q=1-p$. Let
$x_0,x_1,\ldots,x_d$ be a simple path, with $d\ge2$. Suppose that
\[
 \mu(x_j\leftrightarrow x_k)\ge c>0
 \qquad(0\le j<k\le d,\ k-j<d).
\]
For an integer $s\ge1$, suppose that the clusters of $x_0,x_d$ are
distinct and both have at least $s$ vertices with probability
$\delta>0$. There is a path edge $e$ such that
\[
 -\log\mu(H_s(e))
 \le C\bigl(\log(1/\delta)+\lceil\log_2d\rceil+1\bigr).
\]
Here $H_s(e)$ is the event that the endpoint clusters of $e$ are
distinct and each has at least $s$ vertices. The constant $C$ depends
only on $c$ and positive lower bounds for $p$ and $q$.
\end{proposition}

Recall that for probability laws $P,R$ on a finite space,
\[
 D(P\Vert R)=\sum_x P(x)\log\frac{P(x)}{R(x)}
\]
is their relative entropy, with zero summands omitted. We will construct
a law $Q_\eta$ supported on $H_s(e)$ for a fixed path edge $e$, with
relative entropy of order $\log(1/\delta)+\log d$ from $\mu$.
Since any law supported on an event of probability $u$ has relative
entropy at least $\log(1/u)$, this will prove the proposition.

We construct $Q_\eta$ by opening contacts along two chains of clusters.
A reverse procedure that guesses and closes these contacts will account
for the possible multiplicity of the opening map.

An exact incidence exposure of a cluster records its vertex set and the
states of every edge with at least one endpoint in that set, including
internal edges skipped by a lazy search. All edges wholly outside the
exposed set retain their product law. We will repeatedly use this
observation for unions of distinct clusters.

\subsection{Contact weights and normalized exposures}

We first record the normalization that makes the construction possible.
Suppose a union $O$ of exact clusters has been exposed, a deterministic
target set $\mathcal A$ is contained in $O$, and $w\notin O$.
Run the deterministic lazy depth-first search of Lemma~\ref{lem:dfs}
from $w$ towards $\mathcal A$. For every query $f$, let $h$ be the
conditional probability of $w\leftrightarrow\mathcal A$ under the
\emph{original product law $\mu$, conditioned only on this search's
own query prefix}. Write $h_0,h_1$ for the two continuations and
$d_f=h_1-h_0\ge0$. Thus $h=ph_1+qh_0$.

A contact is a queried edge from the newly explored cluster into $O$.
Every contact is closed in the exposed configuration. Give it weight
\[
 a_f=pd_f,
 \qquad A=\sum_{\text{queried contacts }f}a_f,
 \qquad h_* =\mu(w\leftrightarrow\mathcal A).
\]
The numbers $d_f$ are not recomputed under the law conditioned on $O$.
This distinction is essential: under that conditioned law a contact is
forced closed, and the ordinary conditional probability drops by
$h-h_0=pd_f$. Every noncontact query is a fresh Bernoulli bit, so its
conditional expected change in $h$ is zero. The search fails and its
final prefix certifies failure even under $\mu$, since it has queried
all edges leaving the cluster of $w$. Telescoping therefore gives
\begin{equation}\label{eq:contact-normalization}
 \EE_{\mathrm{ext}}A=h_*.
\end{equation}
Here the expectation samples only the product exterior of the fixed
incidence exposure.

Conditional on the fixed incidence exposure and any query prefix, the
same telescoping gives expected remaining contact weight at most one.
The potential $h$ still uses only the search's own prefix. Since each
contact weight is at most one,
expansion of the square yields
\begin{equation}\label{eq:contact-second-moment}
 \EE_{\mathrm{ext}}A^2
 \le \EE_{\mathrm{ext}}\sum_f a_f
       +2\EE_{\mathrm{ext}}\sum_f a_f
 \le3.
\end{equation}
In the middle expression the second sum bounds each contact weight
multiplied by the conditional expected sum of subsequent weights.
Consequently, when $h_*>0$, the rule
\begin{equation}\label{eq:marked-contact-law}
 \frac{a_f}{h_*}\,\mu_{\mathrm{ext}}(d\omega)
 \quad\text{on exterior configurations and a chosen contact }f
\end{equation}
defines a probability measure, with counting measure on the contact
coordinate.

After this choice, expose the exact incidence configuration of the
new cluster $K=C_w$. The chosen contact, its query prefix, and its
weight are determined by the old exposure and the incidence states of
$K$. Thus the factor in \eqref{eq:marked-contact-law} imposes no
condition on edges wholly outside $O\cup K$. Conditional on the new
exposure and chosen contact, those edges still have product law.
This proves both normalization and preservation of the exterior law
at each successive exposure.

\subsection{Binary interpolation and its description}

Let $J=\lceil\log_2d\rceil$. Start with the product law conditioned
on the endpoint event in Proposition~\ref{prop:interpolation}. Expose
the two endpoint clusters, calling them the left and right bases, and
assign every vertex in each base its respective side.
Maintain an interval of the path whose endpoints have opposite sides.
At a step $i$, take its midpoint $w_i$, using the lower index in a tie,
and set
\[
 \mathcal A_i=\{x_0,x_d,w_1,\ldots,w_{i-1}\}.
\]
Every target anchor belongs to a previously exposed cluster. If $w_i$
already lies in their union $O$, make no new exposure and give it its
cluster's side; call this a hold. Otherwise apply
\eqref{eq:marked-contact-law} with $w=w_i$ and
$\mathcal A=\mathcal A_i$, expose $K_i=C_{w_i}$, and record its
chosen contact $f_i$. The cluster containing the head of $f_i$ is its
parent, and $K_i$ inherits that parent's side. Leave $f_i$ closed.
Retain the subinterval joining $w_i$ to the endpoint of opposite side.

The new interval has length at most the ceiling of half the old length,
so after $\ell\le J$ steps its endpoints are adjacent. At each
nonhold step, one current endpoint is a target anchor at path-index
distance strictly less than $d$ from $w_i$. Therefore
\begin{equation}\label{eq:interpolation-root-bound}
 h_i:=\mu(w_i\leftrightarrow\mathcal A_i)\ge c.
\end{equation}
The preceding normalization constructs a joint law of a common final
configuration $\Omega$ and its chosen contacts. Equivalently, one
successively multiplies the conditional exterior law by the normalized
factors \eqref{eq:marked-contact-law}; each factor is measurable in the
new frozen incidence data, so later exterior sampling preserves all
earlier factors.

The parent relations form a forest rooted at the two bases. Call a
nonbase cluster active if it lies on the ancestry path from the owner
of one terminal endpoint to its base. The active clusters on each side
form one chain. Record a descriptor consisting of the number of steps,
the side and hold bit at each step, and the active indices. There are
at most $\sum_{j=0}^J8^j\le8^{J+1}$ descriptors. Condition on one of
maximum probability and denote the resulting joint law by $Q$.
All midpoints, target sets, nonhold indices $I$, active indices $I_*$,
and the terminal edge are now deterministic. Moreover, the parent of
an active index is the preceding active index on its side, or its base
if there is none. Indeed, the active indices on that side are exactly
one ancestry chain, ordered by creation time. Thus every active parent
has a deterministic anchor.

\begin{figure}[tbp]
\centering
\begin{tikzpicture}[x=1cm,y=1cm,
 cluster/.style={draw=blue!55!black,fill=blue!5,ellipse,
 minimum width=1.12cm,minimum height=.82cm,inner sep=5pt},
 contact/.style={draw=blue!65!black,line width=1pt},
 closed/.style={draw=black!45,dashed,line width=.6pt}]
 \node[cluster] (L) at (-5,0) {$L$};
 \node[cluster] (K1) at (-2.8,0) {$K_1$};
 \node[cluster] (K3) at (-.6,0) {$K_3$};
 \node[cluster] (K4) at (1.9,0) {$K_4$};
 \node[cluster] (R) at (4.1,0) {$R$};
 \node[cluster,draw=black!40,fill=black!3] (K2) at (-2.8,-1.55) {$K_2$};
 \draw[contact] (L)--node[above=3pt] {$f_1$}(K1);
 \draw[contact] (K1)--node[above=3pt] {$f_3$}(K3);
 \draw[contact] (K4)--node[above=3pt] {$f_4$}(R);
 \draw[closed] (K3)--node[above=3pt] {$e$}(K4);
 \draw[closed] (K1)--node[right=4pt] {$f_2$}(K2);
 \node[below=5pt] at (L.south) {left base};
 \node[below=5pt] at (R.south) {right base};
 \node[below=5pt,black!65] at (K2.south) {inactive};
\end{tikzpicture}
\caption{The two ancestry chains at the end of interpolation.
Here the terminal path edge $e$ joins $K_3$ to $K_4$.
Only $f_1,f_3,f_4$ are opened. Each terminal cluster joins its own
large base; the two sides remain distinct. The inactive contact
$f_2$ stays closed. Vertices inside each cluster are suppressed.}
\label{fig:interpolation}
\end{figure}


Relative to $\mu$ on configurations and counting measure on the
contacts indexed by $I$, the joint density satisfies
\begin{equation}\label{eq:interpolation-density}
 \frac{dQ}{d(\mu\otimes\mathrm{count})}
 \le e^L\prod_{i\in I}a_i,
 \qquad a_i=pd_{f_i},
 \qquad
 L=\log(1/\delta)+(J+1)\log8+J\log(1/c).
\end{equation}
The three costs are the endpoint conditioning, descriptor conditioning,
and the normalizers \eqref{eq:interpolation-root-bound}.

\subsection{The two chains and their inverse}

Open only the active contacts in $\Omega$, and call the resulting
configuration $\eta$. Each terminal endpoint is now joined along its
active ancestry chain to a base cluster of size at least $s$. The two
resulting clusters remain distinct: all original cluster boundaries were
closed, and the only newly opened edges join clusters within one of the
two same-side chains. Thus $Q_\eta$ is supported on $H_s(e)$ for the
terminal edge fixed by the descriptor.

We next give a reverse procedure and compute its probability of recovering
the chosen contacts. This computation will identify the quantities that
the remaining entropy estimates must control.

For the fixed schedule, expose all anchor clusters in order on an
ordinary product configuration. Let $A_i$ be the sum of the contact
weights at step $i$, with $A_i=1$ when $w_i$ is already exposed. These
quantities depend only on the configuration and schedule, not on earlier
contact choices: the exposed union is exactly the union of the clusters
of the preceding anchors. Write
\[
 X=\sum_{i\in I}\log^+A_i,
 \qquad \log^+t=\log\max\{1,t\}.
\]

For each active $i$, let $t_i$ be a standalone deterministic transcript
that exposes the two exact clusters at $w_i$ and its parent anchor,
querying all their incidence edges and caching duplicate queries.
For an ordinary pair transcript $t$, call a contact $f$ admissible if
the cluster at $w_i$ avoids $\mathcal A_i$ and its depth-first search
queries $f$ closed into the distinct parent cluster. Every pair selected
under $Q$ is admissible.

Open just $f$ in an admissible configuration. It becomes the sole open
exit from the original cluster of $w_i$, and the parent contains a target
anchor. Hence $f$ is a queried actual open pivot. Let $S_i(t,f)$ be the
sum of the weights $qd$ of queried actual open pivots through $f$ in this
modified search, including $f$ itself. This number is determined by
$(t,f)$. The prefix through $f$ uses only incidences of the original
cluster at $w_i$; an earlier queried open edge is an actual pivot exactly
when deleting it separates $w_i$ from the tail of $f$ within that cluster.
The exact incidence transcript records all the internal edges needed for
this test.

Construct a reference probability $R$ on output configurations and marks
as follows. First sample $\eta$ under $\mu$. On a working copy, process
the active indices in descending order. At index $i$, run the
$\mathcal A_i$ depth-first search from $w_i$. For each queried actual open
pivot of weight $b_f=qd_f$, let $S_f$ be the cumulative weight through it,
and give that pivot guess probability
\begin{equation}\label{eq:interpolation-inverse-guess}
 \pi_i(f)=\frac{b_f}{(1+S_f)^2}.
\end{equation}
The sum of these probabilities is at most
$\int_0^\infty(1+t)^{-2}\,dt=1$, by comparison with the integral on
each cumulative-weight interval. Assign the remaining probability to
a dummy guess. Close the guessed pivot and continue, using arbitrary
dummy defaults after an invalid history. After all active guesses, on
the recovered working copy guess each inactive contact independently
with probabilities $a_f/A_i$, again using a dummy default if necessary.
The configuration marginal of $R$ is $\mu$.

We verify that the true inverse guesses have precisely the weights just
defined. At active step $i$, all later active contacts have already been
closed. Every earlier active contact has both endpoints in original
clusters created before $K_i$, since a contact joins its new cluster to
an older parent. Thus none of those earlier contacts touches $K_i$.
Its only open boundary edge in the current working copy is $f_i$, and
its parent still contains the deterministic target anchor. The search
prefix through $f_i$, its original-product conditional differences
$d_f$, and the criterion for earlier actual pivots are therefore unchanged
from the configuration with only $f_i$ opened. In particular, the guess
$f_i$ is available with cumulative weight exactly $S_i$. Holds add target
anchors but no new cluster or switched edge, so they do not affect this
argument.

The active contacts are distinct: different nonhold indices create
distinct original clusters, and each contact joins its new cluster to
an older one. Consequently the map
$(\Omega,(f_i))\mapsto(\eta,(f_i))$ is injective. Along the true inverse
history, \eqref{eq:interpolation-density} bounds the ratio of the joint
output density to $R$ by
\[
 e^L\prod_{i\in I_*}
 \frac{(q/p)a_i}{qd_{f_i}/(1+S_i)^2}
 \prod_{i\in I\setminus I_*}A_i
 =e^L\prod_{i\in I_*}(1+S_i)^2
       \prod_{i\in I\setminus I_*}A_i.
\]
Here $q/p$ is the original-to-opened Bernoulli likelihood ratio for each
active contact. It cancels because the forward weight is $a_i=pd_{f_i}$
and the inverse guess has numerator $qd_{f_i}$.
Taking logarithms and discarding the marks by data processing gives
\begin{equation}\label{eq:interpolation-cost-to-bound}
 D(Q_\eta\Vert\mu)
 \le L+2\EE_Q\sum_{i\in I_*}\log(1+S_i)+\EE_Q X.
\end{equation}
It remains to bound the two expectations on the right by a constant
times $L+J$. We do this under the closed-contact law $Q$, before opening
any edge.

\subsection{Entropy of the sampled configuration}

We use the same formula for $D(P\Vert R)$ when $R$ is a nonnegative finite
measure rather than a probability; in that case it need not be
nonnegative. For $P\ll R$, Jensen's inequality gives
\begin{equation}\label{eq:interpolation-entropy-inequality}
 \EE_P g\le D(P\Vert R)+\log\sum R e^g.
\end{equation}
Write $Q_\Omega$ for the configuration marginal of $Q$ and put
$D_\Omega=D(Q_\Omega\Vert\mu)$.

By \eqref{eq:contact-second-moment},
\[
 \EE_\mu[\max(1,A_i)^2\mid\text{previous exact exposures}]\le4.
\]
Each $A_i$ is known after the next exact exposure. Successive conditional
expectation therefore gives
$\EE_\mu e^{2X}\le4^J$. Summing the density bound
\eqref{eq:interpolation-density} over the possible contact choices,
which only enlarges the sum if the descriptor restrictions are dropped,
gives
\[
 D_\Omega\le L+\EE_Q X.
\]
Applying \eqref{eq:interpolation-entropy-inequality} to $2X$ gives
$2\EE_Q X\le D_\Omega+J\log4$. Hence
\begin{equation}\label{eq:interpolation-configuration-entropy}
 D_\Omega\le2L+J\log4,
 \qquad \EE_Q X\le L+J\log4.
\end{equation}

\subsection{Entropy of the inverse guesses}

We now estimate the sum involving $S_i$. The pair transcripts $t_i$
contain just the cluster data needed to determine each marked weight.
On the support of $Q$, every cluster occurs in at most two such pairs,
and an edge is incident to at most two clusters. Each bit is therefore
read by at most four pair algorithms. Lemma~\ref{lem:entropy} gives
\begin{equation}\label{eq:interpolation-pair-entropy}
 \sum_{i\in I_*}D(Q_{t_i}\Vert\mu_{t_i})\le C_pD_\Omega,
\end{equation}
with $C_p$ uniform when $p,q$ are bounded below. The multiplicity
condition is needed only on the support of $Q$.

For each active $i$, define a finite measure on marked pair transcripts by
\[
 \nu_i(t,f)=\mu_{t_i}(t)\,pd_f
 \quad\text{for admissible }(t,f),
\]
and zero otherwise. This reference measure combines the original pair
transcript law with the contact weight used in the forward sampling.
It need not have mass one. Its usefulness is the following bound, which
converts the cost of an inverse guess to the pivotal second moment.

The map that opens $f$ and records $f$ is injective on full marked
configurations. The original-to-modified product probability ratio is
$q/p$, so multiplying the original weight $pd_f$ gives the modified
weight $qd_f$. In any modified configuration, the admissible inverse
marks are a subset of its queried actual open pivots. Their cumulative
weights are at most the total $W$ from Lemma~\ref{lem:dfs}. Therefore
\begin{equation}\label{eq:interpolation-contact-measure}
 \sum_{t,f}\nu_i(t,f)(1+S_i(t,f))
 \le\EE_\mu[W(1+W)]\le4.
\end{equation}

We next control the entropy of the chosen marks, not just that of their
pair transcripts. Write $\mathcal H(\cdot\mid\cdot)$ for conditional
Shannon entropy. Taking logarithms in
\eqref{eq:interpolation-density} and averaging yields
\begin{equation}\label{eq:interpolation-joint-contact-cost}
 D_\Omega+\sum_{i\in I}\EE_Q\log(1/a_i)
       -\mathcal H((f_i)_{i\in I}\mid\Omega)\le L.
\end{equation}
Since $t_i$ is determined by $\Omega$,
\[
 \mathcal H((f_i)_{i\in I}\mid\Omega)
 \le\sum_{i\in I_*}\mathcal H(f_i\mid t_i)
    +\sum_{i\in I\setminus I_*}\mathcal H(f_i\mid\Omega).
\]
For an inactive index, nonnegativity of conditional relative entropy
to the probabilities $a_f/A_i$ gives
\[
 \EE_Q\log(1/a_i)-\mathcal H(f_i\mid\Omega)
 \ge-\EE_Q\log A_i.
\]
Combining these inequalities with
\eqref{eq:interpolation-configuration-entropy} and
\eqref{eq:interpolation-pair-entropy}, we obtain
\begin{align*}
 \sum_{i\in I_*}D(Q_{t_i,f_i}\Vert\nu_i)
 &=\sum_{i\in I_*}\left[
 D(Q_{t_i}\Vert\mu_{t_i})+
 \EE_Q\log(1/a_i)-\mathcal H(f_i\mid t_i)\right]\\
 &\le C_pD_\Omega+L+\EE_Q X
 \le C'_p(L+J).
\end{align*}
All chosen weights are positive, and all spaces are finite, so these
expressions are finite on the support of $Q$.
Apply \eqref{eq:interpolation-entropy-inequality} with
$g=\log(1+S_i)$ and use
\eqref{eq:interpolation-contact-measure}. This proves
\begin{equation}\label{eq:interpolation-inverse-cost}
 \EE_Q\sum_{i\in I_*}\log(1+S_i)\le C''_p(L+J).
\end{equation}

Substituting \eqref{eq:interpolation-configuration-entropy} and
\eqref{eq:interpolation-inverse-cost} into
\eqref{eq:interpolation-cost-to-bound} yields
\begin{equation}\label{eq:interpolation-output-entropy}
 D(Q_\eta\Vert\mu)\le C'''_p(L+J).
\end{equation}
We already proved that $Q_\eta$ is supported on $H_s(e)$. If
$u=\mu(H_s(e))$, conditioning $\mu$ on this event gives
$D(Q_\eta\Vert\mu)\ge\log(1/u)$.
Together with the definition of $L$, the last display proves
Proposition~\ref{prop:interpolation}.
